% Generated draft: requires mathematical and source review.
\documentclass[11pt]{article}
\usepackage[margin=25mm]{geometry}
\usepackage{fontspec}
\setmainfont{Latin Modern Roman}
\usepackage{amsmath,amssymb,mathtools,unicode-math}
\setmathfont{Latin Modern Math}
\usepackage{enumitem,needspace,xurl,hyperref}
\setlength{\parindent}{0pt}
\setlength{\parskip}{5pt}
\setlength{\emergencystretch}{3em}
\newcommand{\N}{\mathbb{N}}
\newcommand{\Z}{\mathbb{Z}}
\newcommand{\Q}{\mathbb{Q}}
\newcommand{\R}{\mathbb{R}}
\newcommand{\C}{\mathbb{C}}
\newcommand{\F}{\mathbb{F}}
\newcommand{\A}{\mathbb{A}}
\newcommand{\PP}{\mathbb{P}}
\newcommand{\E}{\mathbb{E}}
\newcommand{\eps}{\varepsilon}
\newcommand{\dd}{\,\mathrm{d}}
\begin{document}
\section*{UnsolvedMath ID: 20000342}
\textbf{Source problem code:} catalogue entry
\subsection{Definitions and mathematical statement}
The recovered Problem 4.8 asks to study the geometry of Hacking's moduli space of plane curves and applications to families of smooth plane curves, with analogous questions for configuration spaces and marked del Pezzo surfaces. The degree, stability coefficients, marking data, and intended geometric properties are not specified; this is a research program rather than one quantified conjecture. The verification field says the appended beginning of Section 5 in the raw extraction is outside the problem environment.

The report treats only a generic quartic boundary case, in a characteristic-zero interpretation. Consider a very general smooth pair
\[
(\PP(1,1,4),D),\qquad D:\ w^2=f_8(x,y),
\]
where \(f_8\) is a smooth binary octic. It asserts that the full graded-automorphism calculation gives \(\operatorname{Aut}(\PP(1,1,4),D)=\mu_2\), generated by the hyperelliptic deck involution. Using Hacking's smoothness/dimension theorem and the two-Veronese \(\Q\)-Gorenstein smoothing
\[
x_0x_2-x_1^2=t x_3,
\]
the report proposes a completed local stack chart
\[
\left[\operatorname{Spf}k[[a_1,\ldots,a_5,t]]/\mu_2\right],
\]
where the involution fixes the five binary-octic directions \(a_i\) and sends \(t\) to \(-t\). Thus the asserted coarse completed normal coordinate is
\[
q=t^2,
\]
with coarse ring \(k[[a_1,\ldots,a_5,q]]\). The smoothing coordinates and their weighted ambient model, the stable-pair coefficients, and the exact local-moduli theorem hypotheses are not supplied in the summary.

Over a strictly henselian DVR \(V\) with uniformizer \(\pi\) and \(2\) invertible, a nonzero coarse normal arc \(q=u\pi^m\), \(u\in V^\times\), is asserted to lift to this stack precisely when \(m\) is even. The strict henselian and residue-characteristic conditions are needed to take square roots of units and trivialize the relevant \'etale \(\mu_2\)-torsor. In particular a transverse arc, \(m=1\), requires a quadratic ramified base change, or a square-root stack base. An arc identically in \(q=0\) is separate from this finite-valuation criterion. The candidate claims minimal ramification index two for an ordinary transverse coarse arc.

The report gives an independent genus-three representation check:
\[
H^0(D,2K_D)
=H^0(\PP^1,\mathcal O(4))_+\oplus H^0(\PP^1,\mathcal O)_-,
\]
with dimensions \(5+1\), matching the invariant boundary directions and anti-invariant smoothing direction. Its quartic local calculation does not settle the full geometry or the original configuration/marked-del-Pezzo analogues. Characteristic \(2\), where the deck group and square-root analysis change, is not silently included.
\subsection{Short English statement}
The source asks broadly about Hacking's plane-curve moduli and related spaces. The report proposes a generic quartic boundary chart where the involution reverses the smoothing parameter, making the coarse normal coordinate its square and forcing quadratic ramification for transverse lifts.
\subsection{Sources}
\begin{itemize}
\item AIM workshop ``Compact moduli spaces and birational geometry,'' Problem 4.8: \url{https://aimath.org/WWN/birational/birational.pdf}. No PDF was accessed here.
\item Local canonical registry \texttt{aim-algebraic-geometry-notes.json}, \texttt{notes[341]}, tag \texttt{problem}; permanent UnsolvedMath ID \(20000342\), source code \texttt{catalogue entry}. The supplied verification claims the original problem ends before Section 5 and corrects the extracted contamination. Original remarks and literature field are empty.
\item The raw extraction additionally includes adjacent Section 5.1, ``Birational geometry,'' with references [FG03], [GKM02], [FP], described there as an incomplete list. No full entries are supplied, and these are neighboring-section context rather than references establishing the quartic candidate.
\item Supplied research attempt 1 invokes Hacking's smoothness/dimension theorem and proposes the local deformation and lifting result. No full authoritative bibliographic references or complete stability setup are supplied. The report is machine generated, not peer reviewed, and requires independent verification before citation.
\end{itemize}
\end{document}
